\documentclass[a4paper]{article}
\usepackage{graphicx}
\usepackage{amsmath, amssymb}
\usepackage{csquotes}
\usepackage[
backend=biber,
style=alphabetic,
sorting=ynt
]{biblatex}
\addbibresource{bibliografi.bib}

% Egna kommandon:
% Delar-operator:
\newcommand{\divides}{\mid}
\newcommand{\notdivides}{\nmid}
% Största gemensamma delare
\renewcommand{\gcd}[1]{(#1)}
% Kardinalitet
\newcommand{\card}[1]{\# #1}
% Heltalsmängd:
\newcommand{\A}{\mathcal{A}}
% Primtalsmängd:
\renewcommand{\P}{\mathcal{P}}
% Sållad mängd:
\renewcommand{\S}[3]{S(\mathcal{#1}, \mathcal{#2}, #3)}

%\usepackage{epsfig}
\usepackage{floatflt} %för inkapslade bilder.
\usepackage{epstopdf}
\usepackage{fancyhdr}
\usepackage[top=3cm, bottom=3cm,inner=3cm, outer=3cm]{geometry}	
\setlength{\headheight}{61pt}
%\addtolength{\textwidth}{20mm}
%\addtolength{\textheight}{30mm}
%\addtolength{\textheight}{10mm}
%\addtolength{\headheight}{-10mm}
%\addtolength{\oddsidemargin}{-10mm}
\usepackage{eso-pic}								% Create cover page background
\newcommand{\backgroundpic}[3]{
	\put(#1,#2){
	\parbox[b][\paperheight]{\paperwidth}{
	\centering
	\includegraphics[width=\paperwidth,height=\paperheight,keepaspectratio]{#3}}}}


% Här nedan kommer kommandon där du skall göra val
%


\newcommand{\ARBETE}{ %välj arbetsbenämning här, notera att båda kan användas samtidigt
 \newline \noindent Examensarbete för kandidatexamen i matematik vid Göteborgs universitet % om någon i gruppen läser på GU
% \medskip
% \newline \noindent Kandidatarbete inom civilingenjörsutbildningen vid Chalmers % om någon läser på Chalmers
 \medskip
}


\newcommand{\titel}{Titel} %Skriv in projektets/rapportens titel här
\newcommand{\undertitel}{Eventuell undertitel} %Skriv in ev undertitel här, eller lämna tomt
\newcommand{\engtitel}{Engelsk titel} %Skriv in engelsk översättning av projektets/rapportens titel här


\newcommand{\namn}{ %Skriv in medlemmarnas namn i bokstavsordning.
    Kawthar Abdalla \\
    Nils Alexandersson \\
    Erik Dagobert \\
    Coën Lorcan Olofsson
}

\newcommand{\examina}{ % här skall gruppens medlemmar skrivas in vid önskad examen. Aktivera aktuella examina, inte kurskoder. Vid fyra med samma examen blir det snyggast om man gör en tabell.
%  \newline \noindent \begin{tabular}{ll}
%  förnamn efternamn 1 &
%  förnamn efternamn 2  \\
%  förnamn efternamn 3 &
%  förnamn efternamn 4
% \end{tabular}
% Vid fem eller sex görs tabellen med tre kolumner
% Vid flera examina skall \bigskip aktiveras utom för den sista.
%
%
%%%%%%%%%%%%%%%%%% Kurs MMG900 %%%%%%%%%%%%%%%%
% \newline \noindent {\it Examensarbete för kandidatexamen i matematik vid Göteborgs universitet} \smallskip
% \newline \noindent förnamn efternamn 1
% \quad förnamn efternamn 2
% \quad förnamn efternamn 3
% \quad förnamn efternamn 4
% \bigskip
%%%%%%%%%%%%%%%%%% Kurs MMG910  %%%%%%%%%%%%%%%%
\newline \noindent{ \it  Examensarbete för kandidatexamen i matematik inom Matematikprogrammet vid Göteborgs universitet} \smallskip
% \newline \noindent förnamn efternamn 1
% \quad förnamn efternamn 2
% \quad förnamn efternamn 3
% \quad förnamn efternamn 4
% \bigskip
%
\newline \noindent \begin{tabular}{@{}ll}
    Kawthar Abdalla &
    Nils Alexandersson \\
    Erik Dagobert &
    Coën Lorcan Olofsson
\end{tabular}
\bigskip
%%%%%%%%%%%%%%%%%% Kurs MMG920  %%%%%%%%%%%%%%%%%
%  \newline \noindent {\it Examensarbete för kandidatexamen i matematik inom Matematikprogrammet, inriktning Tillämpad matematik, vid Göteborgs universitet} \smallskip
%  \newline \noindent förnamn efternamn 1
%  \quad förnamn efternamn 2
%  \quad förnamn efternamn 3
%  \quad förnamn efternamn 4
%  \bigskip
%
%%%%%%%%%%%%%%%%%%% Kurs MSG900  %%%%%%%%%%%%%%%%%
% \newline \noindent {\it Examensarbete för kandidatexamen i matematisk statistik vid Göteborgs universitet} \smallskip
% \newline \noindent förnamn efternamn 1
% \quad förnamn efternamn 2
% \quad förnamn efternamn 3
% \quad förnamn efternamn 4
% \bigskip
%
%%%%%%%%%%%%%%%%%%% Kurs MSG910  %%%%%%%%%%%%%%%%%
%\newline \noindent {\it Examensarbete för kandidatexamen i matematisk statistik inom Matematikprogrammet vid Göteborgs universitet} \smallskip
%
%\newline \noindent förnamn efternamn 1
%\quad förnamn efternamn 2
%\quad förnamn efternamn 3
%\quad förnamn efternamn 4
% \bigskip
%
%%%%%%%%%%%%%%%%%%% Kurs MVEX01  %%%%%%%%%%%%%%%%%
% \newline \noindent {\it Kandidatarbete i matematik inom civilingenjörsprogrammet Automation och mekatronik vid Chalmers} \smallskip
% \newline \noindent förnamn efternamn 1
% \quad förnamn efternamn 2
% \quad förnamn efternamn 3
% \quad förnamn efternamn 4
 % \bigskip
%
% \newline \noindent {\it Kandidatarbete i matematik inom civilingenjörsprogrammet Datateknik vid Chalmers} \smallskip
% \newline \noindent förnamn efternamn 1
% \quad förnamn efternamn 2
% \quad förnamn efternamn 3
% \quad förnamn efternamn 4
% \bigskip
%
% \newline \noindent {\it Kandidatarbete i matematik inom civilingenjörsprogrammet Maskinteknik vid Chalmers} \smallskip
%
% \newline \noindent förnamn efternamn 1
% \quad förnamn efternamn 2
% \quad förnamn efternamn 3
% \quad förnamn efternamn 4
% \bigskip
%
%   \newline \noindent {\it Kandidatarbete i matematik inom civilingenjörsprogrammet Teknisk fysik vid Chalmers} \smallskip
%   \newline \noindent förnamn efternamn 1
%   \quad förnamn efternamn 2
%   \quad förnamn efternamn 3
%   \quad förnamn efternamn 4
%   \bigskip
%
%   \newline \noindent {\it Kandidatarbete i matematik inom civilingenjörsprogrammet Teknisk matematik vid Chalmers} \smallskip
%   \newline \noindent förnamn efternamn 1
%   \quad förnamn efternamn 2
%   \quad förnamn efternamn 3
%   \quad förnamn efternamn 4
%   \bigskip
%
 \bigskip

}


\newcommand{\handledare}{% om samtliga handledare är från MV lämnas fältet inst tomt
Lucile Devin \\
& Anders Södergren
%  namn 1& inst \\
%  & namn 2 & inst \\
%  & namn 3 & inst \\
}



% Här nedan kommer kommandon som du inte skall ändra, de ger utformningen av rapporten.

\newcommand{\skribenter} {\begin{tabular}{l} \namn \end{tabular}}

%%%%%%%%%%%%%%%%% Här börjar utformningen av omslaget %%%%%%%%%%%%%%%%
\newcommand{\omslag}{
%\newgeometry{top=3cm, bottom=3cm,left=2.25 cm, right=2.25cm}	% Temporarily change margins	
\newgeometry{top=3cm, bottom=4cm,left=2 cm,right=1cm}	% Temporarily change margins	
%\thispagestyle{empty}
% \addtolength{\textheight}{-20mm}
\pagestyle{fancy}
\pagenumbering{gobble}
\fancyhead[C]{\includegraphics[width=170mm]{Chalmers_GU_svart-eps-converted-to.pdf}\\}
\addtolength{\voffset}{0.3cm}
\renewcommand{\headrulewidth}{1pt}


\parbox{17cm}{
\vspace{60mm}

\noindent{\Huge \titel}
\bigskip

\noindent {\Large \undertitel}
\bigskip

\noindent{\huge \engtitel}

\noindent\hspace*{-1 ex}{\Large \it
\ARBETE
}
\vspace{20mm}

\noindent\hspace*{-1 ex}\parbox{80mm}{\noindent {\huge \skribenter}}}

\renewcommand{\footrulewidth}{1pt}

\fancyfoot[L]{\vspace{0.1mm}\large Institutionen för Matematiska vetenskaper\\
CHALMERS TEKNISKA HÖGSKOLA\\
GÖTEBORGS UNIVERSITET\\
Göteborg, Sverige 2020 }


\newpage
\thispagestyle{empty}
\mbox{}

\newpage}
%%%%%%%%%%%%%% Utformning av omslag slut %%%%%%%%%%%%%%%%%%%%

%%%%%%%%%%%%%%%%%%%% Här börjar utformning av titelsidor %%%%%%%%%%%%%%%%
\newcommand{\titelsidor}{
\newgeometry{top=3cm, bottom=5cm,left=3 cm,right=3cm}

\thispagestyle{fancy}
\fancyhf{}
\renewcommand{\headrulewidth}{0pt}


\mbox{}
\vspace{50mm}

\noindent {\LARGE \titel}\bigskip\bigskip

\noindent {\large \undertitel}
\vspace{30mm}


\noindent {\large \examina}


\vfill
\hspace{-5.8 ex} \begin{tabular}[t]{lll}
Handledare:& \handledare & \end{tabular}


\renewcommand{\footrulewidth}{0pt}

\fancyfoot[L]{\vspace{0.1mm}\large Institutionen för Matematiska vetenskaper\\
CHALMERS TEKNISKA HÖGSKOLA\\
GÖTEBORGS UNIVERSITET\\
Göteborg, Sverige 2020 }




\newpage
\thispagestyle{fancy}
\fancyhf{}
\newgeometry{top=3cm, bottom=3cm,left=3 cm,right=3cm}
\mbox{}
\vfill

\setcounter{page}{0}
\newpage }



\usepackage[T1]{fontenc}                % För svenska bokstäver
\usepackage[swedish,english]{babel}             % För svensk avstavning och svenska rubriker (t ex "Innehållsförteckning")


\begin{document}
\selectlanguage{swedish}
\omslag

\titelsidor
\thispagestyle{empty}
%\setlength{\textheight}{240mm}
%\addtolength{\topmargin}{-50mm}
\newgeometry{top=3cm, bottom=3cm,left=3 cm,right=3cm}

\section*{Förord}

I förordet skall det anges vilka delar som skall tillskrivas respektive författare. Det skall också anges att en loggbok förts över de enskilda medverkandes prestationer under arbetet. Med loggbok menas här gruppens dagbok och de individuella tidsloggarna.

\newpage

\section*{Populärvetenskaplig presentation}

Här skriver man populärvetenskaplig presentation

\newpage
\begin{abstract}

\end{abstract}

\selectlanguage{english}
\begin{abstract}

\end{abstract}

\newpage
\selectlanguage{swedish}
\pagestyle{plain}
\tableofcontents                % Innehållsförteckning

                % Tabellförteckning

\newpage
\pagenumbering{arabic}
\section{Inledning}
Den som någon gång har funderat kring primtal, kanske har provat att ta en lista med naturliga tal upp till någon gräns och börjat försöka hitta de primtalen som finns. 
Efter en liten stund kanske man börjar lägga märke till mönster som uppträder; såsom att det finns vissa par av primtal som bara har ett tal mellan sig. 
Man kanske också börjar stryka igenom alla multipler av primtal för att ha något system för att uteslutta sammansatta tal.
Så småningom kanske man inser att för att hitta alla primtal upp till ett visst tal behöver man bara alla primtal mindre än eller lika med kvadratroten av det talet.

Idén bakom denna process, att hitta primtal i en lista av naturliga tal, har funnits sedan antikens grekland med en algoritm som har tillskrivits den grekiska polyhistorn Eratosthenes (ca. 276 - 194 f.v.t.). Algoritmen har följande struktur;
\begin{enumerate}
    \item Börja en lista med talet 2 och lista alla naturliga tal upp till någon gräns.
    \item Ringa in 2 och stryk över alla andra tal som är delbara med 2.
    \item Ringa in nästa tal som inte är struket och stryk alla andra tal delbara med det nya inringade talet.
    \item Upprepa steg 3 tills varje tal på listan är struket eller inringat. 
\end{enumerate}
När algoritmen är avslutad så har varje primtal i listan blivit inringat och alla andra tal har strukits. 
Eratosthenes algoritm lade grunden för ämnet sållteori; ett område inom matematiken som försöker uppskatta storkleken på så kallade \textit{siktade mängder}. 

En siktad mängd är en mängd där alla element är heltal och har någon gemensam egenskap som till exempel en mängd som består av endast primtal eller mängden av alla tal i en aritmetisk talföjld.
Den största fördelen med den grundläggande sållteorin är att metoderna är relativt elementära och flexibla, speciellt jämfört med andra metoder inom analytisk talteori. 
Det krävs inga idéer från komplex analys som till exempel Dirichlets L-funktioner eller serier för att ha användning av de enklare sållen och om man kan formulera mängderna på ett korrekt sätt, går det att effektivt tillämpa metoderna på ett stort antal mängder. 
Effektiviteten innebär att trots sin enkla konstruktion, kan dessa matematiska såll fortfarande ge starka resultat, även om bättre resultat kan hittas genom att använda mer raffinerade analytiska metoder.
Exempelvis gäller det att antalet primtal mindre än \textit{x} har det asymptotiska beteendet av \(x/\log(x)\) vilket ges i primtalssatsen. 
Detta icke trivialaresultat kan nästan bevisas utan användning av zeta-funktioner eller då man utnyttjar Selbergs såll, dock så får man bara en övre gräns av \(x/\log(x)\) istället för det asymptotiska beteendet. 
Mer avancerade sållmetoder har givit svar på frågor närliggande till primtalstvillingsförmodan i \cite{chen2Prime}, och att det finns oändligt många par av primtal med maximalt 600 tal emellan sig \cite{mayBound}.

Vår rapport fokuserar på små såll av både kombinatorisk och viktad form. 
Först, att sållen+ är små innebär det att deras sållning använder ett få antal restklasser modulo primtal.
 Näst, att ett såll är kombinatoriskt innebär att den använder sig av inklusion-exklusionsprincipen för att uppskatta mängdens storlek på ett lämpligt sätt. 
Slutligen, att ett såll är viktat innebär att någon viktfunktion används för att uppskatta storleken. 
Sållen vi håller oss till är Eratosthenes generaliserade såll, Bruns såll, och Selbergs såll. 
För varje såll ger vi en kortfattad härledning till dess formulering, där vi lyfter fram de centrala idéerna för dess konstruktion, och en förklaring till hur sållet används med exempel.
En diskussion angående de sållens feletermerna följer sista teoretiska tillämpningen.
Efteråt redovisar och analyserar vi en implementering i kod av Eratosthenes ursprungliga algoritm, vilket följer metoden som beskrivs i \cite{HaraldSieve}.
Med hjälp av de primtal som genereras av algoritmen undersökar vi slutligen vissa egenskaper hos primtalen såsom deras fördelning, fördelningen av primtalstvillingar, och frekvensen av olika stora avstånd mellan två på varandra följande primtal (primtalsgap).
Dock, innan vi börjar redovisa något av sållen går vi igenom några förberedelser och förklarar det allmänna sållproblemet.

\section{Förkunskaper}
% Skrivet av Erik

%\subsection{O-notation}

\subsection{Det allmänna sållproblemet}
% Kanske kan koppla till inklusions-exklusionsprincipen. 
Det generella fallet i sållteori är att vi har en mängd heltal \(\A\), en mängd primtal \(\P\) samt delmängder \(\A_p, \forall p \in \P\) och är intresserade av att sålla fram delmängden \(\mathcal{S}(\A, \P) := \A \setminus \cup_{p \in \P} \A_p\) eller, vanligare, kardinaliteten av dito. För \(d\), en kvadratfri produkt (dvs. utan delare \(p^\alpha\) för \(\alpha > 1\)) av element \(p \in \P\) definierar vi \(\A_d = \cap_{p \divides d} \A_p\). Om \(P(z) = \prod_{\substack{p\in \mathcal{P} \\ p < z}} p\) så betecknar vi kardinaliteten av den sållade mängden som
\begin{align*}
    \S{A}{P}{z} = \card{(\A \setminus \cup_{p \divides P(z)} \A_p)} .
\end{align*}


%\subsection{O-notation}
% Intresserade av en övre asymptotisk gräns --> förklara Ordonotation. 

\subsection{Möbiusfunktionen}
En användbar funktion i sållteori är möbiusfunktionen,
\begin{equation*}
    \mu(n) = 
    \begin{cases}
        1, & \text{om}\ n \text{ är ett kvadratfritt, naturligt tal med jämnt antal primdelare}\\
        -1, & \text{om}\ n \text{ är ett kvadratfritt, naturligt tal med udda antal primdelare}\\
        0, & \text{om}\ n \text{ inte är kvadratfri}
    \end{cases}
\end{equation*}
som bland annat förenklar inklusion-exklusionsprincipen så att, för scenariot ovan,
\begin{align*}
    \S{A}{P}{z} = \sum_{\substack{d \divides P(z) \\ d \leq z}} \card{\A_d} . %Mobius?
\end{align*} % sida 72

Två andra viktiga egenskaper hos möbiusfunktionen är den så kallade fundamentala egenskapen,
\begin{equation*}
    \sum_{d \divides n} \mu(d) =
    \begin{cases}
        1, & \text{if}\ n = 1 \\
        0, & \text{if}\ n > 0
    \end{cases}
\end{equation*}
och Möbius inverteringsformel som säger
\begin{equation*}
    f(n) = \sum_{d \divides n} g(d) \implies g(n) = \sum_{d \divides n} \mu(d) f(n/d)
\end{equation*}
om \(f, g : \mathbb{N} \to \mathbb{C}\)


%\subsection{Abels summationsformel}

\section{Eratostenes generella såll}


\subsection{Eratosthenes–Legendre Såll}
\hspace{0.3cm} Den gregiska matematikern, astronomen och poeten Eratosthenes från Cypern (can276 f.v.t – ca 194 f.v.t), var en av de tidigaste matematikerna som arbetade med primtal. I en av hans stora presentationer inom matematik introducerade han ”Eratosthenes såll”, vilket är ett effektivt sätt att identifiera primtal upp till en viss gräns, som senare blev grunden för sållteorin.
År 1808 introducerade den franska matematikern Adrien-Marie Legendre (1752-1833) det som kom att kallas Legendres Såll i sin bok "Théorie des nombres", som är en formulering av Eratosthenes-metoden, som heter "Eratosthenes-Legendre såll" eftersom den bygger på Eratosthenes algoritm. Hur har denna såll startat den moderna sållen och vilka upptäckter gjordes med hjälp av Eratosthenes-Legendre såll?\\



 Legendre formaliserade matematiken bakom Eratosthenes algoritmen genom att använda den inklusiv-exklusivprincipen för att beräkna antal elementer i en given mängd av naturliga tal  $ A\subset Z_{+} $ som inte har några primtalsfaktorer nedan $ z $. Med $  \#A=X $, tar vi alla $ x $ tal, och vi substraherar antalet multiplar av 2, även vi substraherar antalet multiplar av 3, samma sak med multiplar av 5, sen lägger vi till antalet av multiplar $ 2\times3 $ (eftersom multipler av 6 substraherades två gånger, en gång med multiplarna av 2 och den andra med dem av 3); och vi fortsätter på detta sätt genom alla värden på  $ d $ som delar produkten av primtal upp till  $ z $. \\
 

Låt $ P_z $ vara produkten av all primtal i $ \P $ mindre än $ z $, vi har
\[P_z:=\prod_{p<z}p\]
Låt $ S(\A, P) $ beteckna antal heltal i $ \A $ som är relativt prima till $ P_z $, således
\[S(\A, P)=\vert \lbrace n\in \mathbb{N}:1\leq n\leq x, (n,P_z)=1 \rbrace\vert\]


Legendres identiteten uttrycks av: 

\begin{theorem}[Legendres identitet \cite{Terence}] kan uttryckas genom:
\[S(\A, P)=\sum_{d\divides P_z}\mu(d)\sum_{\substack{n\leq x;\\d\divides n}}1. \]
% Where the $ S(\A, P(z)) $ calculate the number of natural numbers $ n\leq x $ in $ A $ coprime to $ P(z) $.
\end{theorem}
Vi kommer att använda denna sats för att komma till en allmän övre gräns för $ S(\A, P) $\\

Vi har
\begin{align*}
S(\A, P) & =\sum_{d\divides P_z}\mu(d)\Bigl\lfloor \frac{x}{d}\Bigr\rfloor\\
 &  =\sum_{d\divides P_z}\mu(d) \left( \frac{x}{d}+ \Bigl\lfloor \frac{x}{d} \Bigr\rfloor - \frac{x}{d} \right)\\
 & =\sum_{d\divides P_z}\mu(d)\frac{x}{d}+\sum_{d\divides P_z}O(1)\\
 & =\sum_{d\divides P_z}\mu(d)\frac{x}{d}+ O\left(\sum_{d\divides P_z}1\right)
\end{align*}

Så, vi kan hitta en uppskattning för den inre summan i Legendres identitet som följande:
\[\sum_{n\leq x;d\divides n}1=\frac{x}{d}+O(1); \]
Vi vet att antalet åtskilda faktorer som $ P(z) $ har är $ 2^{\pi(z)} $ vi härleda
\[S(\A, P)=\sum_{d\divides P_z}\mu(d)\frac{x}{d}+ O(2^{\pi(z)}) \]
För varje $ p $ i $ \P $, vi betecknar $ \omega(p) $ antalet av utvalda restklasser mod $ p $ i mängden $ \lbrace \omega_{1,p},...,\omega_{\omega_{p},p}\rbrace $ för varje primtal i $ \P $, vi definerar $ W(z) $ som
\[W(z)=\prod_{\substack{p\in \P\\p<z}}\left( 1-\frac{\omega(p)}{p}\right)\]
Genom att faktorisera den första termen av summan 
\[\sum_{d\divides P_z}\mu(d)\frac{x}{d}= x\prod_{p< z}\left( 1-\frac{1}{p} \right) \]
och med ersättning med den nya faktoriserad termen, vi får
\[S(\A, P)=\left( \prod_{p< z}\left( 1-\frac{1}{p} \right)\right) x +O(2^{\pi(z)}) \]




\subsection{Eratosthenes Generella Såll}
Efter att vi såg den grundläggande sållen, realiserar vi att den inte ger en nyttig uppskattning av $ S(\A, P) $ på grund av dess stora feltermen. Den grundläggande metoden av Eratosthenes-Legendre såll kan generaliseras, vi kan anpassa den för att den kan användas med andra talföljder av siffror som kvadratfria heltal. En kvadratfri heltal är en heltal som är en produkt av åtskilda primtal. Dess mer generella såll kommer att vara mycket bättre och effektivare med de kvadratfria tal i grunden eftersom de är tätare än primtal, detta kommer att hjälpa oss att hitta intressanta gränser för dessa kvadratfria siffror.

Den generalla såll kan användas för att demonstrera många intressanta resultater. I denna del kommer vi att introducera Eratosthenes generella såll och dess bevis.

Efter att vi sett Legendres identitet i sats 3.1 och förutom definitionerna i förgående avsnittet, så kan vi skriva om sållen och formulera den som följande:


Låt $ P(z) $ vara produkten av alla primtal i $ \P $ mindre än $ z $, d.v.s
\[P(z)=\prod_{\substack{p \in \P\\ p \leq z}} p \]

Vi definerar $\A_{p}$ som en mängd av elementer av $ \A $ på ett sätt där det tillhör åtminstone en av $\omega(p)$ utvalda restklasser mod p.

Låt $ d $ vara kvadratfri som endast är en produkt av primfaktorer i $ \P $, vi sätter $\omega(d)=\prod_{p \divides d} \omega(p)$

Så vi har
\[\omega(d)=\prod_{p\divides d}\omega(p)\]

Låt $ R_{d} $ representera feltermen i uppskattningen av antal elementerna av $ \A_{d} $ 


Låt $ \kappa $ vara siktdensitetkonstanten, det är ett viktat medelvärdet av antalet restklasser som siktas ut av varje primtal.


Antar att $ \vert R_{d}\vert=O(\omega(d))$ and $ \exists\kappa\geq 0 $ så att
\begin{equation}
\label{equationEL3}
\sum_{p\divides P(z)}\frac{\omega(p)\log(p)}{p}\leq \kappa \log(z)+O(1)
\end{equation}  
Med partiell summering får vi
\[\sum_{p\divides P(z)}\frac{\omega(p)}{p}\leq\kappa\log\log z+O(1)\]
   Vi har de följande Lemman som ska användas för att bevisa Eratosthenes generella såll satsen.
\begin{lemma}
     Antar \ref{equationEL3} har vi
$$
\sum_{d<t, d \divides P(z)} \omega(d)=O\left(t(\log z)^{\kappa} \exp \left(-\frac{\log t}{\log z}\right)\right)
$$
där den stora- $O$ gräns är för $z \rightarrow \infty$ och konstanten bero på $P, R_{p}, \kappa$.
\end{lemma}
\begin{lemma}
    Fixa $C>0 .$ Antar (1) har vi
$$
\sum_{d>C x, d \divides P(z)} \frac{\omega(d)}{d}=O\left((\log z)^{\kappa+1} \exp \left(-\frac{\log x}{\log z}\right)\right)
$$
där den stora- $O$ gräns är för $z \rightarrow \infty$ och konstanten bero på $P, R_{p}, \kappa, C .$
 
\end{lemma}

Med de förgående antagendena, Eratosthenes generella såll är formulerad som följande\cite{Dalton}
\begin{theorem}[Eratosthenes Generella Såll]\hfill

För en given $ R_{d} $ antar vi att $ X $ existerar, som verifierar
\[\vert\A_{d}\vert =\frac{\omega(d)}{d}X +R_{d} \]

Låt oss anta $ \exists \kappa\geq 0 $ och $ \vert R_{d}\vert = O(\omega(d))$ på ett sätt så att

\[\sum_{p\divides P(z)}\frac{\omega(p)\log(p)}{p}\leq \kappa \log(z)+O(1)\]


Med (3) och genom att anta $ \exists y\in \mathbb{R}^{+} $ som verifierar $ \forall d>y, \card{\A_d}=0 $.

får vi
\begin{equation}
\label{equationEL4}
\S\A\P z= XW(z)+O\left( \left( X+\frac{y}{\log(z)} \right)(\log(z))^{\kappa+1}\exp \left( -\frac{\log(y)}{\log(z)} \right) \right)
\end{equation}
\end{theorem}
\begin{proof}
Med de hypoteser som används för Eratosthenes generella sållteorin och den inklusion-exklusionprincipen, får vi
\begin{align*}
\S\A\P z &=\sum_{\substack{d\divides P(z)\\d\leq y}}\mu(d)\#\A_{d} =\sum_{\substack{d\divides P(z)\\d\leq y}}\mu(d)\frac{X\omega(d)}{d}+O(\sum_{\substack{d\leq y,\\ d \divides P(z)}} \omega(d)).
\end{align*}
Så genom att ersätta med de förgående två Lemman, får vi (4)
\end{proof}


\subsection{Applikationer och begränsningar}

En huvudapplikation, är att finna en övre gräns för $ \pi(x) $, mycket gjordes för att förbättra resultatet. Här är några:
\begin{proposition}
    Genom att använda Eratosthenes–Legendre såll, ser vi att den övre gränsen av antalet primtal upp till $ x $ är given av
\[S(\A, P)\ll\frac{x}{\log(\log(x))}\]
\end{proposition}
\begin{proof}
Vi har från sats 3.1
\[S(\A, P)=x \prod_{p\leq z}\left( 1-\frac{1}{p} \right) +O(2^{\pi(z)})\]
För att få en övre gräns för $ S(\A, P) $, vi får en nedre gräns för den reciproka värdet av första termen, 

Genom att använda de geometriska serier, får vi
\[\prod_{p\leq z}\left( 1-\frac{1}{p} \right)^{-1}=\prod_{p\leq z}\sum_{r=0}^{\infty}\frac{1}{p^{r}}\]
Så har vi
\[\prod_{p\leq z}\sum_{r=0}^{\infty}\frac{1}{p^{r}}>\sum_{n<z}\frac{1}{n}\]
För att termen på höger sidan är en delmängd av dem som är på vänster sidan.\\

Genom att jämföra med $ \int_{1}^{z}\frac{1}{x}dx $, får vi
\[\sum_{n<z}\frac{1}{n}>\log(z)\]
Således
\[S(\A, P)<\frac{x}{\log(z)}+O(2^{\pi(z)})\]
För att kontrolera den andra termen (felterm), låt $ z=\log(x)$. Vi vet att $ \pi(z)<z $ så $ 2^{\pi(z)}<2^{z}=2^{\log(x)}=x^{log(2)} $, således får vi
\[S(\A, P)\ll\frac{x}{\log(\log(x))}\]
\end{proof}
Efteråt, en förbättring gjordes, vilket gav en bättre övre gräns för antal primtal upp till $ x $, den ges av nästa förslag
\begin{proposition}
En annan övre gräns för antal primtal upp till x är given av
\[S(\A, P)\ll\frac{x}{\log(x)}(\log\log x).\]
\end{proposition}

En annan viktig tillämpning av Eratosthenes–Legendre såll är att hitta en övre gräns för antalet primstvillning upp till $ x $.
\begin{theorem}[Twin prime upper bound\cite{Aliram}]\hfill

Antal primtal $ p $ with $ p\leq x $ och $ p+2 $ är också en primtal (d.v.s primtalstvillning), är given av
\[\S\A\P z \ll\frac{x}{\log^{2}(x)}(\log\log x)^{2}\]
\end{theorem}
\begin{proof}
Låt $ \A $ beteckna vilken mängs av positiva heltal, $ \P $ betecknar vilken mängd av primtal. Låt ett reelt tal $ z $ such as $ z=z(x) $. Vi skiljer residualklasser 0 och -2 modulo $ p $ för varje primtal $ p<z $ . Given att mängden $ \A_p $ är nollmängd för $ p>x+2 $, och med användning av Eratosthenes Generella Såll (sats 3.4) för $ \kappa=2 $ får vi
\[\S\A\P z = xW(z)+O \left( x(\log z)^{3} \exp \left( -\frac{\log x}{\log z}\right) \right)\]
med
\[W(z)=\prod_{p<z}\left( 1-\frac{2}{p}\right)\]
Vi kan hitta den övre gränsen för $ W(z) $ som följande
\[W(z)=\prod_{p<z}\left( 1-\frac{2}{p}\right)\leq\exp \left( -\sum_{p<z}\frac{2}{p} \right)\ll (\log z)^{-2}\]
Om vi sätter värdet av $ z $ som
\[\log z = \log x/A\log\log x\] med $ A $ en stor positiv konstant
\[\S\A\P z \ll \left( \frac{\log x}{A\log\log x} \right)^{-2}=\left( \frac{A\log\log x}{\log x} \right)^{2}\]
Därför, sammanfattar vi
\[\S\A\P z \ll\frac{x}{\log^{2}(x)}(\log\log x)^{2}\]
\end{proof}
Notera att denna övre gränsen är inte den bästa. 
 


Eratosthenes–Legendre Såll, har många tillämplingar i såll och talteorin. Efter dess presentation, många andra sållmetoder definierades som förlängningar av Eratosthenes–Legendre såll som Brun's såll.


%\begin{thebibliography}{9}

%\bibitem{Terence} 
%Terence Tao. 
%Notes 7: Sieving and expanders, 1 March, 2012,\\
%\url{https://terrytao.wordpress.com/2012/03/01/254b-notes-7-sieving-and-expanders/}

%\bibitem{Dalton} 
%Jack Robert Dalton. 
%An Exposition of Selberg’s Sieve,  Master thesis,  University of Vermont, May, 2017, pages: 32-33.
%\\\url{https://scholarworks.uvm.edu/cgi/viewcontent.cgi?article=1719&context=graddis}

%\bibitem{Aliram}
%Alina Carmen Cojocaru, M. Ram Murty.
%An Introduction to Sieve Methods and Their Applications, Cambridge University Press, London Mathematical Society, 2005, page 73.


%\end{thebibliography}



\section{Bruns såll}
% Skrivet av Nils
\begin{comment}
När den norska matematikern Viggo Brun (1881-1978) år 1915 publicerade artikeln \textit{Über das Goldbachsche Gesetz und die Anzahl der Primzahlpaare}, lades grunden till modern sållteori \cite{ViggoBrun}. 
I artikeln presenterade Brun ett såll baserat på Eratosthenes idéer, som numera kallas Bruns såll och fyra år senare använde han det för att bevisa att summan av reciproker $\frac{1}{p_1}+\frac{1}{p_2}+...$ av primtalstvillingar $p_1,p_2,...$ konvergerar \cite{ViggoBrun}.
I detta avsnitt ges en kortfattad beskrivning av Bruns såll och ett exempel taget ur \cite{cojocarumurty} på hur sållet kan tillämpas för att ge ett resultat som angränsar till primtalstvillingsförmodan.
\end{comment}

När den norska matematikern Viggo Brun (1881-1978) år 1915 publicerade artikeln \textit{Über das Goldbachsche Gesetz und die Anzahl der Primzahlpaare}, lades grunden till modern sållteori \cite{ViggoBrun}. 
I artikeln presenterade Brun ett såll baserat på Eratosthenes idéer, som numera kallas Bruns såll och fyra år senare använde han det för att bevisa Bruns sats, vilken finns beskriven i avsnitt \ref{eratosthenes.tillämpning} ovan.
I detta avsnittet ges en kortfattad beskrivning av Bruns såll och ett exempel taget ur \cite{cojocarumurty} på hur sållet kan tillämpas för att frambringa ett resultat som angränsar till primtalstvillingsförmodan.


\subsection{Beskrivning av sållet}
Brun visade att för varje funktion $g$ som uppfyller $g(1)=1$,
så gäller det att
\begin{equation*}
    \S{A}{P}{z} 
    = \sum_{d \mid P(z)} \mu(d) g(d) \card\A_d 
    - f_g,
\end{equation*} % f_g = \sum_{d \mid P(z)} \sum_{\substack{p \mid P(z) \\ p<q(d)}} \mu(d) (g(d)-g(pd)) S(\A_{pd},\P^{(pd)},p).
där $f_g$ är en funktion av $\A,\P$ och $z$, definierad med hjälp av $g$. 
Därefter konstruerade Brun två funktioner $g_U$ och $g_L$ så att $f_{g_U}\geq 0$ och $f_{g_L}\leq 0$ alltid gäller. 
Således kan en undre och övre begränsning till $\S{A}{P}{z}$ erhållas, nämligen
\begin{equation}\label{brun.eq.gugl}
    \sum_{d \mid P(z)} \mu(d) g_L(d) \card\A_d 
    \leq \S{A}{P}{z} 
    \leq \sum_{d \mid P(z)} \mu(d) g_U(d) \card\A_d.
\end{equation}
Enligt \cite{cojocarumurty} var detta Bruns innoverande idé.
Genom att sedan skriva $\card\A_d=\frac{\omega(d)}{d}X + R_d$ kunde han, efter mycket möda, hitta en approximation och en felterm för begränsningarna i (\ref{brun.eq.gugl}) och på så vis formulera Bruns såll. Följande sats är hämtad ur \cite[Kap. 6.2]{cojocarumurty}.


\begin{theorem}[Bruns såll] \label{brun.thm.brun}
Med notation från avsnitt \ref{sallproblemet}, låt $b$ vara ett positivt heltal och $\lambda$,
ett reellt tal som uppfyller $0<\lambda e^{1+\lambda}<1$.
Antag att $\left|R_d\right|\leq\omega(d)\leq dA_1$ för varje kvadratfritt tal $d$ med faktorer ur $\P$ och någon konstant $A_1<1$,
samt att det finns konstanter $A_2\geq 1$ och $\kappa>0$ som uppfyller
\begin{align*}
    \sum_{w\leq p<z} \frac{\omega(p)\log p}{p} \leq \kappa\log(z/w) + A_2,\quad \text{för}\ 2\leq w\leq z.
\end{align*}
Då begränsas $\S{A}{P}{z}$ ovanifrån av
\begin{align}
    & XW(z)\left(1 + \lambda c_1\exp\left(c_2\frac{2b+3}{\lambda\log z}\right)\right) + O\left(z^{c_3}\right), \label{brun.eq.upper}
    \intertext{och underifrån av}
    & XW(z)\left(1 - c_1\exp\left(c_2\frac{2b+2}{\lambda\log z}\right)\right) + O\left(z^{c_3 - 1}\right). \label{brun.eq.lower}
\end{align}
Konstanterna $c_1$, $c_2$ och $c_3$ definieras som
\begin{equation*}
    c_1 := \frac{ 2\lambda^{2b}e^{2\lambda} }{ 1 - \lambda^2e^{2+2\lambda} }, \quad
    c_2 := \frac{A_2}{2}\biggl(1+\frac{\kappa+\frac{A_2}{\log 2}}{1+A_1}\biggr)\ \text{och} \quad
    c_3 := 2b + \frac{2.01}{e^{2\lambda/\kappa} - 1}.
\end{equation*}
\end{theorem}
Precis som i avsnitt \ref{Eratosthenes} betraktar vi här $\kappa$ som dimensionen av sållet.
Notera att konstanten $b$ kan väljas relativt fritt och kan bidra till att hålla nere feltermen.
Nedan följer ett exempel på hur detta kan göras, där en tillämpning av satsen presenteras.


\subsection{En tillämpning av Bruns såll}
Låt oss uttrycka en modifierad version av primtalstvillingsförmodan på följande vis:\\
\textit{Det finns oändligt många par av heltal $(n,n+2)$ där båda talen har högst $r$ primtalsfaktorer.}\\
För valet $r=1$ erhålls exakt primtalstvillingsfömodan, vilken står obevisad.
Däremot går det att med hjälp av Bruns såll bevisa påståendet för $r=7$, vi följer samma tillvägagångssätt som i \cite[Kap. 6.2]{cojocarumurty}.
Definiera mängden
\begin{equation*}
    \mathcal{T} := \{\textit{$n\leq x$: $n$ och $n+2$ har som mest $7$ primtalsfaktorer}\}.
\end{equation*}
Målet är att hitta en undre begränsning till $\card{\mathcal{T}}$ som divergerar då $x\to\infty$.
Låt $\P$ vara mängden av alla primtal och låt
\begin{equation*}
    \A := \{n(n+2): n\leq x\}.
\end{equation*}
Med hjälp av (\ref{brun.eq.lower}), vill vi nu uppskatta storleken på $\A$ efter bortsållning av tal delbara med primtal $p<z$ för något $z$.
Vi skiljer på två typer av restklasser; noll och nollskild modulo $p$, och har därmed att
\begin{equation*}
    \omega(2)=1\ \text{ och }\ \omega(p)=2,\ \text{för}\ p>2.
\end{equation*}
Detta ger direkt att $\frac{\omega(p)}{p}\leq\frac{2}{3}$ med likhet då $p=3$, låt därför $A_1=2/3$.
För att hitta ett passande $\kappa$, använder vi oss av (\ref{Thm.1.4.4}) som tillsammans med $\omega(p)\leq2$ implicerar
\begin{equation*}
    \sum_{w\leq p<z} \frac{\omega(p)\log p}{p} \leq 2\log\left(z/w\right) + O(1),
\end{equation*}
varpå vi ser att $\kappa=2$ uppfyller kriterierna och att $A_2$ existerar.
Något explicit värde på $A_2$ är inte nödvändigt, vilket kan ses genom att betrakta exponenten i ($\ref{brun.eq.lower}$). Taylorutveckling ger
\begin{align}\label{brun.example.exp}
    \exp\left(c_2\frac{2b+2}{\lambda\log z}\right) = 1 + c_2\frac{2b+2}{\lambda}(\log z)^{-1} + \left(c_2\frac{2b+2}{\lambda}\right)^2(\log z)^{-2} + ...
\end{align}
Fixera $b,\lambda$ och $A_2$ så att även $c_2$ är fixerat, då följer det att (\ref{brun.example.exp}) begränsas av $1+O((\log z)^{-1})$, då $z\to\infty$.
Vidare gäller det att 
\begin{align*}
    W(z) = \frac{1}{2}\exp\Bigl(\sum_{2<p<z} \log(1-2/p) \Bigr) 
    \gg \frac{1}{2}\exp\Bigl(\sum_{2<p<z} \frac{-2}{p} \Bigr).
\end{align*}
Partiell summation av (\ref{Thm.1.4.4}) ger att $\sum_{p\leq n}\frac{1}{p} = \log\log z + O(1)$, så
\begin{align*}
    W(z) \gg \exp\Bigl(-2\log\log z\Bigr)
    = (\log z)^{-2}.
\end{align*}
Genom tillämpning av sats \ref{brun.thm.brun} erhålls
\begin{equation} \label{brun.example.sap}
    \S{A}{P}{z} \gg \frac{x}{(\log z)^2}\bigl(1 - c_1\bigr) + O\left(z^{c_3 - 1}\right).
\end{equation}
Nu vill vi välja $b$ och $\lambda$, så att $c_1<1$ och $c_3$ är litet. 
Detta uppfylls av $b=1$ samt $\lambda=0.101$ som medför att $c_3-1<8$.

I fallet att $n$ eller $n+2$ har mer än 7 primtalsfaktorer,
så måste $p\mid n(n+2)$ för något primtal $p < n^{1/8}$.
Således erhålls $\card\mathcal{T}$ ur $\S{A}{P}{z}$ genom att sålla för primtal upp till $z = x^{1/u}$, där $u<8$.
Insättning av detta i (\ref{brun.example.sap}) ger
\begin{equation*}
    \frac{u^2x}{(\log x)^2}\bigl(1 - c_1\bigr) + O\left(z^{\frac{c_3 - 1}{u}}\right) \gg 
    \frac{x}{(\log x)^2} + O\left(z^{\frac{c_3 - 1}{u}}\right).
\end{equation*}
Genom att ställa det ytterligare kravet $u>c_3-1$, har vi att feltermen ovan försvinner när $x$ blir stort och vi får slutligen
\begin{equation} \label{brun.example.bound}
    \card\mathcal{T} \gg \frac{x}{(\log x)^2}.
\end{equation}
Notera här att vi inte bara har bevisat påståendet \textit{det finns oändligt många par av heltal $(n,n+2)$ där båda talen har högst 7 primtalsfaktorer},
uttrycket (\ref{brun.example.bound}) ger oss även en undre gräns för det asymptotiska beteendet hos mängden ifråga.
Härnäst riktas fokus mot det tredje och sista sållet i denna text.






%\section{Metod och genomförande}
%% Skrivet av Nils

Denna rapport är framför allt en litteraturstudie, med stor förankring i \textit{An Introduction to Sieve Methods and Their Applications} av Alina Carmen Cojocaru \& M. Ram Murty (2005). Rapporten ämnar \textbf{att introducera matematikstudenter på sen kandidatnivå till sållteori.} Vi presenterar tre stycken såll som har haft stor teoretisk och historisk betydelse för ämnet och exemplifierar dessa i programmeringsproblem. 

I teoridelen formuleras de mest centrala satserna och idéerna men bevisen har i vissa fall kortats ned och i andra fall helt överlåtits till appendix, detta i syfte att göra texten mer lättillgänglig för läsaren.



%\begin{thebibliography}{FFF}
%\bibitem[UR]{rapp} Utformning av rapporter och kandidatarbetens skriftliga presentation för Civilingenjörsprogrammen vid Chalmers tekniska högskola. 2008. Göteborg: Chalmers Tekniska Högskola

%\end{thebibliography}
\medskip

\printbibliography

\end{document}                 % The input file ends with this command.
